\documentclass[10pt,a4paper]{amsart}
\usepackage[margin=19mm]{geometry}
\usepackage{amsmath,amssymb,mathtools}
\usepackage[scr=rsfs,cal=euler]{mathalfa}
\usepackage{xfrac}
\usepackage[unicode,hidelinks,bookmarks=false]{hyperref}
\newcommand{\ie}{i.e.\ }
\newcommand{\cf}{cf.\ }
\newcommand{\resp}{respectively\ }
\newcommand{\Subsection}{\subsection}
\providecommand{\nobreakdash}{\nobreak\hbox{-}}
\setlength{\emergencystretch}{2em}
\multlinegap0pt
\usepackage{amssymb}
\usepackage{mathtools}
\usepackage{xcolor}
\usepackage{tikz}
\newtheorem{lemma}{Lemma}
\newtheorem{corollary}{Corollary}
\newtheorem{theorem}{Theorem}
\newcommand{\bN}{\mathbb N}
\newcommand{\cZ}{\mathcal Z}
\title{A measurable equivariant Weierstrass theorem}
\author{Konstantin Slutsky, Mikhail Sodin, Aron Wennman}
\date{Source: arXiv:2607.05542v1 (CC BY 4.0)}
\begin{document}
\maketitle

\noindent\emph{Source and licence.} A measurable equivariant Weierstrass theorem. Version arXiv:2607.05542v1 (CC BY 4.0), distributed by arXiv under the Creative Commons Attribution 4.0 International licence, http://creativecommons.org/licenses/by/4.0/, as recorded in the arXiv licence metadata for this exact version and shown on the abstract page https://arxiv.org/abs/2607.05542v1. Full licence notice: \texttt{licenses/arxiv-2607.05542-CC-BY-4.0.txt}.\par\smallskip

\noindent\textit{Editorial scope.} Counted unit: the article's corollary cor:indep\_sets (source0.tex lines 1536--1539). The article's complete original proof of it is delivered with it (source0.tex lines 1541--1573, one \texttt{proof} environment of 33 lines). No mathematics is written, repaired or re-derived: the statement and the proof are the article's own lines, in the article's own notation, and no retained line is substituted or re-flowed.  Retained uncounted as the source-local context the proof needs: lines 1501--1506; lines 2190--2215.  Nothing was omitted from any retained span, so the package carries no omission record.  No retained line contains a citation, so the wrapper carries no bibliography and no bibliography entry.  No retained line contains a figure environment, an image-inclusion command or drawing code, so no image is delivered and the package declares no asset dependency.  Editorial formatting only, in addition: an AMS \texttt{amsart} preamble that restates the article's own declarations and macros used by the retained lines, namely the environments \texttt{lemma}, \texttt{corollary}, \texttt{theorem} on separate counters, together with the standard packages those lines need.  The retained environments print in this document as Lemma 1, Corollary 1, Theorem 1 in order of appearance, whereas the article prints the same items with the numbering of its own full document.  The complete native source of the article as distributed in the arXiv e-print for this exact version is bundled unchanged as \texttt{sources/204549/source0.tex}.
\begin{lemma}[maximal independent sets]\label{lemma:indep_sets_independent_cover}
Let $X$ be a standard Borel space and let $G\subset X\times X$ be a Borel graph. Assume that
$G$ is locally countable. If \(X\) can be
covered by a countable family of Borel independent sets, then each Borel independent set $J\subseteq X$
can be extended to a maximal Borel independent set.
\end{lemma}

\begin{theorem}[ordered finite-section form of the Luzin--Novikov theorem]\label{thm:L-N-ordered}
Let $X$ and $Y$ be standard Borel spaces, let $<$ be a Borel linear order on $Y$, and let
$B\subseteq X\times Y$ be Borel with finite sections
\[
B_x=\{y\in Y\colon (x,y)\in B\}.
\]
Then the counting function
\[
N_B\colon X\to\bN \cup \{0\},\qquad N_B(x)=|B_x|,
\]
is Borel.

Moreover, for every $j\ge1$, there is a Borel map
\[
e_j\colon \{x\in X\colon N_B(x)\ge j\}\to Y
\]
such that, whenever $N_B(x)=p$, the section $B_x$ is listed in increasing order as
\[
B_x=\{e_1(x)<e_2(x)<\ldots<e_p(x)\}.
\]
In particular, for each $p$, the map
\[
x\mapsto (e_1(x),\ldots,e_p(x))
\]
is Borel on the Borel set $\{x\colon N_B(x)=p\}$.
\end{theorem}

\begin{corollary}[maximal independent sets]\label{cor:indep_sets}
Let $X$ be a standard Borel space and let $G\subset X\times X$ be a Borel graph. Assume that
$G$ is locally finite. Then $G$ admits a Borel maximal independent set of vertices.
\end{corollary}

\begin{proof}
  Since all uncountable Borel spaces are isomorphic, and the case of a countable \(X\) is trivial,
  we may assume without loss of generality that \(X = 2^{\bN}\) is the Cantor space.  Consider the
  countable collection $\mathscr Z = (Z_n)_n$ of cylindrical subsets of \(X\), i.e., sets of the
  form
  \[[s] = \{x \in 2^{\bN}\colon  x \text{ extends } s\},\]
  where \(s\) is a finite binary string.  Adjoining the empty set to \(\mathscr Z\), we get a
  countable family of Borel sets that separates points (i.e., for any pair $x\ne x'$ in $X$, there
  exists $\cZ\in\mathscr Z$ such that $x\in \cZ$, while $x'\notin \cZ$) and is closed under finite
  intersections.

  Observe that, for any finite subset $S=\{x_1, \ldots , x_N\}$ of $X$ and any $x\in X\setminus S$,
  there exists $Z\in \mathscr Z$ such that $x\in Z$, while $S\subset X\setminus Z$.  Indeed, for
  each $1\le i \le N$, take $n_i$ so that $x\in Z_{n_i}$, while $x_i\notin Z_{n_i}$, and set
  $Z = \bigcap_{i=1}^N Z_{n_i}$. Since $\mathscr Z$ is closed under finite intersections,
  $Z\in\mathscr Z$.

  For each $x\in X$, the set of its neighbours $\mathcal N(x)=\{y\colon (x, y)\in G\}$ is finite, so we can
  find $Z_n\in\mathscr Z$ that separates $x$ from $\mathcal N(x)$, i.e., such that $x\in Z_n$, while
  $\mathcal N(x)\cap Z_n = \emptyset$. Put
  \[
  I_{n} \stackrel{\rm def}= \{x \in X\colon
n \text{ is minimal such that } x \in Z_{n} \text{ and } \mathcal N(x)\cap Z_n = \emptyset\}.
  \]
  This is an independent subset of $G$, and the family \((I_{n})_{n}\) covers \(X\).  Furthermore,
  these sets are Borel.  Indeed, Theorem~\ref{thm:L-N-ordered} (the ordered finite-section version of
the Luzin--Novikov theorem) gives us Borel uniformizations
  \(g_{k}\colon  X \to X\), \(k \in \bN\), that enumerate neighbours of \(x\). Then
  \[  I_{n} = \{x \in X\colon n \text{ is minimal such that } x \in Z_{n} \text{ and }
    g_{k}(x) \not \in Z_{n} \text{ for all }k \text{ satisfying } g_{k}(x) \ne x\}\]
are Borel, and
Lemma~\ref{lemma:indep_sets_independent_cover} does the job.
\end{proof}

\end{document}
